% !TEX root = ../../../main.tex
\section{Matrix multiplication}
\label{sec:linear-algebra-i:matrix-multiplication}

A matrix with \(r\) rows and \(c\) columns has size \(r\times c\). Write
\(\mathrm M_{r\times c}(\mathbb F)\) for the space of \(r\)-by-\(c\)
matrices with entries in \(\mathbb F\). More generally,
\(\mathrm M_{r\times c}(V)\) denotes the \(r\)-by-\(c\) arrays with entries
in a vector space \(V\).
The notation \(\mathcal M_{r\times c}\) also denotes scalar-valued
matrices of this size when the field is understood.

\subsection{Multiplication by columns}

% Lucas authorized this prerequisite definition on November 9, 2026.
Let \(A=[A_1,\ldots,A_s]\in\mathrm M_{r\times s}(\mathbb F)\).
For a column \(b=(b^1,\ldots,b^s)^T\), multiplication by \(A\) is
defined by the linear combination of its columns
\[
  Ab=\sum_{k=1}^s A_kb^k\in\mathbb F^r.
\]
If \(B=[B_1,\ldots,B_t]\in\mathrm M_{s\times t}(\mathbb F)\),
define its product with \(A\) column by column,
\[
  AB=[AB_1,\ldots,AB_t]\in\mathrm M_{r\times t}(\mathbb F).
\]
Thus the number of columns of \(A\) must equal the number of rows of
\(B\). We write \(I_{r\times r}\) for the identity matrix and \(0\)
for a zero matrix of the required size.

\subsection{Exercises on matrix multiplication}

% !TEX root = ../../hw02.tex
% Source: ma-gy7033fall2026homework2.pdf, p. 1, Problem 1.
% Omit the source's "given in lecture and slides" so the statement is
% self-contained. Correct the typographical "k-the" to "k-th".
\begin{exercise}
  \label{ex:hw02:matrix-product-entries}
  Given a matrix \(M\), let \(M^j_k\) denote the entry in the \(j\)-th
  row and \(k\)-th column of \(M\). Let \(A\in\mathcal M_{r\times s}\)
  and \(B\in\mathcal M_{s\times t}\). Show that the definition of
  matrix multiplication is the same as
  \[
    (AB)^j_l=\sum_{k=1}^s A^j_kB^k_l,
    \qquad 1\le j\le r,\quad 1\le l\le t.
  \]
  \solutionlink{hw02:matrix-product-entries}
\end{exercise}

% !TEX root = ../../hw02.tex
% Source: ma-gy7033fall2026homework2.pdf, p. 1, Problem 2.
% Correct "Let A,B,C are" to "Let A,B,C be" and supply the missing stop.
% Replace "the formula above" with the existing shared exercise reference.
\begin{exercise}
  \label{ex:hw02:matrix-product-associativity}
  Let \(A,B,C\) be scalar-valued matrices such that \((AB)C\) and
  \(A(BC)\) are well defined. Use the formula for matrix multiplication
  in \autoref{ex:hw02:matrix-product-entries} to prove associativity,
  i.e., \((AB)C=A(BC)\).
  \solutionlink{hw02:matrix-product-associativity}
\end{exercise}

% !TEX root = ../../hw02.tex
% Source: ma-gy7033fall2026homework2.pdf, p. 1, Problem 3(a--b).
\begin{exercise}
  \label{ex:hw02:block-shear-inverses}
  For \(X,Y\in\mathcal M_{k\times l}\), let
  \(P,Q\in\mathcal M_{k+l\times k+l}\) be given by
  \[
    P=\begin{bmatrix}I_{k\times k}&X\\0&I_{l\times l}\end{bmatrix}
    \quad\text{and}\quad
    Q=\begin{bmatrix}I_{k\times k}&Y\\0&I_{l\times l}\end{bmatrix}.
  \]
  \begin{enumerate}[(a)][2]
    \item Compute \(PQ\).
    \item Given \(X\), find a matrix \(Y\) such that \(PQ=QP=I\).
  \end{enumerate}
  \solutionlink{hw02:block-shear-inverses}
\end{exercise}

% !TEX root = ../../hw02.tex
% Source: ma-gy7033fall2026homework2.pdf, p. 1, Problem 4(a--b).
\begin{exercise}
  \label{ex:hw02:characterizing-the-identity}
  \begin{enumerate}[(a)][2]
    \item Fix \(s\ge1\). Prove that if \(J\in\mathcal M_{r\times r}\)
      satisfies \(JB=B\) for every \(B\in\mathcal M_{r\times s}\),
      then \(J=I_{r\times r}\).
    \item Show by example that, for every \(r\ge2\), \(JB=B\) for a
      single nonzero \(B\) does not imply \(J=I_{r\times r}\). What
      happens when \(r=1\)?
  \end{enumerate}
  \solutionlink{hw02:characterizing-the-identity}
\end{exercise}

% !TEX root = ../../hw02.tex
% Source: ma-gy7033fall2026homework2.pdf, p. 2, Problem 6(a--f) and hint.
% Source issue: part (f) refers to (d), although universal commutation is
% the condition in (e). Lucas approved the corrected reference November 9, 2026.
\begin{exercise}
  \label{ex:hw02:matrix-commutation}
  Recall that matrix multiplication is not necessarily commutative.
  \begin{enumerate}[(a)][2]
    \item For \(A\in\mathcal M_{r\times s}\) and
      \(B\in\mathcal M_{p\times q}\), determine exactly when \(AB\)
      and \(BA\) are both defined.
    \item Determine when \(AB\) and \(BA\) have the same size.
    \item Given \(r\ge2\), find matrices \(A,B\in\mathcal M_{r\times r}\)
      such that \(AB\ne BA\).
    \item Let \(E^i_j\in\mathcal M_{r\times r}\) be the matrix whose
      entry in row \(i\) and column \(j\) is \(1\) and all of whose
      other entries are \(0\). Prove that
      \(\{E^i_j:1\le i,j\le r\}\) is a basis of
      \(\mathcal M_{r\times r}\) and that
      \[
        E^i_jE^k_l=\delta^k_jE^i_l,
      \]
      where \(\delta^k_j=1\) if \(j=k\) and \(0\) otherwise.
    \item Prove that if \(A\in\mathcal M_{r\times r}\) satisfies
      \(AB=BA\) for every \(B\in\mathcal M_{r\times r}\), then
      \(A=cI_{r\times r}\) for some \(c\in\mathbb F\).
    \item Determine all \(A\in\mathcal M_{r\times r}\) that commute
      with every \(E^i_i\), \(1\le i\le r\). Conclude that, for
      \(r\ge2\), commuting with every diagonal matrix is a strictly
      weaker condition than the one in (e).

  \end{enumerate}
  \emph{Hint.} Compare the entries of \(AE^i_j\) and \(E^i_jA\).
  \solutionlink{hw02:matrix-commutation}
\end{exercise}

% !TEX root = ../../hw02.tex
% Source: ma-gy7033fall2026homework2.pdf, p. 2, Problem 7(a--b).
\begin{exercise}
  \label{ex:hw02:zero-products}
  There are nonzero matrices \(A\) and \(B\) such that \(AB=0\).
  \begin{enumerate}[(a)][2]
    \item Find all \(A\in\mathcal M_{2\times2}\) with \(A^2=0\).
    \item Find \(A,B\in\mathcal M_{2\times2}\) with \(AB=0\) and
      \(BA\ne0\).
  \end{enumerate}
  \solutionlink{hw02:zero-products}
\end{exercise}

% !TEX root = ../../hw02.tex
% Source: ma-gy7033fall2026homework2.pdf, p. 2, Problem 8.
\begin{exercise}
  \label{ex:hw02:upper-triangular-subspace}
  Prove that the set of upper triangular matrices
  \[
    \{A\in\mathcal M_{r\times r}:A^j_k=0\text{ whenever }j>k\}
  \]
  is a subspace of \(\mathcal M_{r\times r}\), compute its dimension,
  and prove that it is closed under matrix multiplication.
  \solutionlink{hw02:upper-triangular-subspace}
\end{exercise}

% !TEX root = ../../hw02.tex
% Source: ma-gy7033fall2026homework2.pdf, p. 2, Problem 9(a--f).
% Omit "be the matrices on the slide Examples of Matrix Multiplication"
% from the shared statement so that its explicit matrices stand alone.
% Source issue: numerical conclusions need field qualifications; none is
% supplied in the PDF. Lucas approved working over R or C November 9, 2026.
\begin{exercise}
  \label{ex:hw02:two-by-two-matrix-identities}
  Let \(\mathbb F=\R\) or \(\C\), and let
  \[
    A=\begin{bmatrix}2&3\\4&-1\end{bmatrix},\qquad
    B=\begin{bmatrix}0&1\\1&-1\end{bmatrix}.
  \]
  \begin{enumerate}[(a)][2]
    \item Compute \(AB-BA\).
    \item Show that \(X\in\mathcal M_{2\times2}\) satisfies \(AX=XA\)
      if and only if \(X=pI+qA\) for some \(p,q\in\mathbb F\).
    \item Show that \(X=0\) is the only \(X\in\mathcal M_{2\times2}\)
      with \(AX=XB\).
    \item Show that \(A^2=A+14I\).
    \item Use this to write \(A^3\) in the form \(pA+qI\).
    \item Use this to find \(N\in\mathcal M_{2\times2}\) such that
      \(AN=NA=I\).
  \end{enumerate}
  \solutionlink{hw02:two-by-two-matrix-identities}
\end{exercise}

